\documentclass[12pt]{article}
\usepackage[margin=0.3in,top=0.25in,bottom=0.4in]{geometry}
\makeatletter
\renewcommand{\@maketitle}{\newpage\null\vskip -3.5em\begin{center}{\LARGE \@title \par}\end{center}\par\vskip 0.2em}
\makeatother
\usepackage{amsmath,amssymb}
\usepackage{booktabs}
\usepackage{listings}
\usepackage{hyperref}
\setlength{\parskip}{0.6em}
\setlength{\parindent}{0pt}
\lstset{basicstyle=\ttfamily\small,columns=fullflexible}

\title{The $\alpha_s$ fix: bottom-quark threshold}
\author{}
\date{}

\begin{document}
\maketitle

\section*{In short}

Our running coupling was 8--10\% too large at the scales of our calculation, because it was
run down from $m_Z$ with four quark flavors all the way, skipping the region above the bottom-quark
mass where five flavors are active. The code now switches from $N_f=5$ to $N_f=4$ at
$m_b=4.8$~GeV. This lowers the diffractive (TMD) cross section by up to 10\% at low $p_{D\perp}$;
the exclusive cross section and $R_{pA}$ do not change.

\section*{1. Where $\alpha_s$ enters}

$\alpha_s$ appears only in the hard factor $\mathcal{H}$ of the diffractive (TMD) cross section,
evaluated at $\mu_R=m_t=\sqrt{p_{D\perp}^2+m^2}$. The exclusive cross section has no $\alpha_s$:
the strength of the strong interaction is contained in the dipole amplitude fitted to HERA data.
In the code, $\alpha_s$ is a prefactor applied in the plotting scripts
(\texttt{src/alphas\_running.py}, used by \texttt{plotting\_scripts/D0.py} and the others).

\section*{2. One-loop running}

At one loop,
\[
  \alpha_s(\mu)=\frac{\alpha_s(\mu_0)}{1+2b_0\,\alpha_s(\mu_0)\ln\left(\dfrac{\mu}{\mu_0}\right)}\,,\qquad
  b_0=\frac{33-2N_f}{12\pi}\,,
\]
which is the same as $1/\alpha_s(\mu)=1/\alpha_s(\mu_0)+2b_0\ln(\mu/\mu_0)$.
How fast $\alpha_s$ grows towards low scales is set by $b_0$, which depends on $N_f$, the number of
quarks lighter than $\mu$. With fewer active quarks $b_0$ is larger, and $\alpha_s$ grows faster.

\section*{3. The problem}

We start from $\alpha_s(m_Z)=0.118$ at $m_Z=91.2$~GeV, where five quarks ($u,d,s,c,b$) are active.
Our scales are $1.5$--$25$~GeV, so going down from $m_Z$ we cross the bottom-quark mass
$m_b\approx4.8$~GeV:
\begin{itemize}
\item between $m_Z$ and $m_b$ the bottom quark is active: $N_f=5$;
\item below $m_b$ only four quarks are active: $N_f=4$.
\end{itemize}
The old code used $N_f=4$ all the way from $m_Z$. Between $m_Z$ and $m_b$ it therefore ran too fast,
and $\alpha_s$ came out too large at our scales.

\section*{4. The fix}

Run with $N_f=5$ from $m_Z$ down to $m_b$, and continue with $N_f=4$ from the value reached at
$m_b$, so that $\alpha_s$ is continuous there (the standard leading-order treatment of a flavor
threshold):
\[
  \alpha_s(\mu)=
  \begin{cases}
    \dfrac{\alpha_s(m_Z)}{1+2b_0(5)\,\alpha_s(m_Z)\ln(\mu/m_Z)}\,, & \mu\geq m_b\,,\\[12pt]
    \dfrac{\alpha_s(m_b)}{1+2b_0(4)\,\alpha_s(m_b)\ln(\mu/m_b)}\,, & \mu<m_b\,,
  \end{cases}
\]
where $\alpha_s(m_b)$ is the value of the first line at $\mu=m_b$. In the code:
\begin{lstlisting}
inv_mb = 1/alphas_mZ + 2*b0(5)*log(mb/mZ)            # 1/alpha_s(mb), Nf = 5 from mZ
alpha_s(mu) = 1/(1/alphas_mZ + 2*b0(5)*log(mu/mZ))   if mu >= mb
            = 1/(inv_mb      + 2*b0(4)*log(mu/mb))   if mu <  mb
\end{lstlisting}
The same change is in \texttt{src/alphas\_running.py} and \texttt{src/alphas\_running.cpp} (the
C++ version is not used by the current programs, but is kept identical).

\section*{5. Size of the change}

\begin{center}
\begin{tabular}{cccc}
\toprule
$\mu$ [GeV] & before ($N_f=4$ throughout) & now (threshold at $m_b$) & before/now \\
\midrule
1.5  & 0.330 & 0.299 & 1.10 \\
3.0  & 0.253 & 0.235 & 1.08 \\
5.0  & 0.216 & 0.203 & 1.07 \\
10   & 0.180 & 0.173 & 1.04 \\
91.2 & 0.118 & 0.118 & 1    \\
\bottomrule
\end{tabular}
\end{center}

\section*{6. Effect on the results}

\begin{itemize}
\item \textbf{Diffractive (TMD) part:} lower by 10\% at low $p_{D\perp}$, falling to about 4\% at
  $p_{D\perp}\approx10$~GeV.
\item \textbf{Exclusive part:} unchanged (no $\alpha_s$).
\item \textbf{$R_{pA}$:} almost unchanged, because the same $\alpha_s$ appears in the Pb result and in
  the proton baseline.
\item \textbf{No reruns needed:} $\alpha_s$ is applied when the plots are made, and all plots have been
  regenerated.
\end{itemize}

\section*{7. Text for the paper}

Instead of ``$N_f=4$ active quark flavors throughout'':
\begin{quote}
The strong coupling $\alpha_s$ is evaluated at one-loop accuracy,
\[
  \alpha_s(\mu_R)=\frac{\alpha_s(\mu_0)}{1+2b_0(N_f)\,\alpha_s(\mu_0)\ln\left(\dfrac{\mu_R}{\mu_0}\right)}\,,
  \qquad b_0(N_f)=\frac{33-2N_f}{12\pi}\,,
\]
starting from $\alpha_s(m_Z)=0.118$ at $\mu_0=m_Z=91.2$~GeV with $N_f=5$ active flavors, and
switching to $N_f=4$ below the bottom-quark threshold $m_b=4.8$~GeV, where $\mu_0=m_b$ and
$\alpha_s(m_b)$ is obtained from the five-flavor running, so that $\alpha_s$ is continuous at $m_b$.
\end{quote}

\end{document}
